
\pagebreak
\begin{center}
\textbf{Laboratory assignment 1: REST API.}
\end{center}

\medskip

\begin{enumerate}

\item Choose on of the following topics for your application:
\begin{itemize}
\item Study Buddy: An app that allows students to connect with each other for studying and sharing resources. It could include features like study groups, file sharing, and chat functionality.

\item Task Tracker: A simple app for managing and tracking tasks and assignments. Students can create to-do lists, set reminders, and mark tasks as completed.

\item Time Management Tool: An app that assists in managing productivity time effectively. It could include features like setting work schedules, creating reminders, and tracking time spent on different tasks.

\item Budget Planner: An app that helps users track their expenses, set budgets, and manage their finances. It could include features like expense categorization, monthly summaries, and notifications for overspending.

\item College Marketplace: A platform for students to buy and sell items within their college community. It could include categories for textbooks, furniture, electronics, and other student necessities.

\item College Event Calendar: A centralized platform for students to discover and keep track of events happening on campus, such as workshops, seminars, and social gatherings. It could include event details, RSVP functionality, and reminders.

\item Grade Calculator: An app that helps students calculate their grades based on different assignments, exams, and weightings. It could also provide insights into how different scores would impact their overall grades.

\item Virtual Lab Assistant: An interactive app that simulates laboratory experiments for science and engineering students. It allows them to perform virtual experiments, collect data, and analyze results.

\item Flashcard Generator: An app that allows students to create digital flashcards for studying different subjects. It could include features like importing images, organizing flashcards into decks, and tracking progress.

\item Internship/Job Board: A platform that aggregates internship and job opportunities specifically targeted at students. It provides a centralized hub for students to find and apply for relevant positions based on their interests and qualifications.

\item Book Exchange Platform: A web app that enables users to trade or lend books with other students in their neighborhood. It could include features like book listings, search filters, and messaging capabilities.

\item Student Feedback System: An app that enables students to provide feedback on courses, professors, and other aspects of their educational experience. It could include anonymous feedback options, rating systems, and comment sections.

\item Virtual Career Fair: An app that brings together employers and students in a virtual setting. It allows companies to showcase job/internship opportunities, conduct interviews, and provide resources for career development.

\item Mental Health Support: An app focused on promoting mental health and well-being. It could include features like guided meditation, mood tracking, stress management techniques, and access to mental health resources.

\item Group Project Manager: An app designed to facilitate group projects by helping students assign tasks, set deadlines, and collaborate effectively. It could include features like file sharing, task tracking, and communication tools.

\item Personalized Book Recommendation: Description: Build a book recommendation app that suggests personalized book lists based on the user's reading preferences, ratings, and genres of interest.

\item Food Waste Reduction: Description: Develop an app to reduce food waste by connecting users with nearby restaurants, cafes, or grocery stores offering surplus food at discounted prices before it goes to waste.

\item Dream Journal with AI Analysis: Description: Create a dream journal app where users can record their dreams. The app uses natural language processing and AI to analyze dream patterns and meanings.

\item Music Collaboration Platform: Description: Develop a platform for musicians and artists to collaborate on music projects virtually, allowing them to record, mix, and produce music together.

\item Local Community Exchange Platform: Description: Create a platform where members of a local community can offer and exchange goods, services, and skills with each other, fostering a sense of community and collaboration.

\item Language Learning Game with AI: Description: Build an interactive language learning game that uses AI to adapt to each user's language proficiency and provide personalized lessons and challenges.

\item Real-Time Collaborative Code Editor: Description: Develop a collaborative code editor that allows multiple users to work on the same codebase simultaneously, promoting real-time collaboration among developers.


\end{itemize} 

\item Generate github repository using a classroom link from moodle course.

\item Choose tech stack for implementation 

\begin{itemize}
\item Node.js with Express

\item Python Django

\item ASP.NET or .Net WebApi
\end{itemize}



\item Build a simple REST API with the following endpoints: 

\begin{itemize}
\item \verb"GET /items" � List all items

\item \verb"POST /items" � Add a new item

\item \verb"GET /items/{id}" � Get details of a specific item

\item \verb"PUT /items/{id}" � Update item info

\item \verb"DELETE /items/{id}" � Remove an item
\end{itemize}

Here instead of items you might use a more descriptive word specific for your task (for example, for budget planner app~--- \verb"budgets")

\item Create a simple frontend (html+css+js /React /Vue.js) that interacts with the REST API endpoints.
    
Example of fetching data from API endpoint:

\begin{lstlisting}
async function getData() {
  const url = "https://your-app-domain//items";
  try {
    const response = await fetch(url);
    if (!response.ok) {
      throw new Error(`Response status: ${response.status}`);
    }

    const json = await response.json();
    update(json); // update your front-end with data from back-end here
  } catch (error) {
    console.error(error.message);
  }
}
\end{lstlisting}

Using different request method (for example POST)

\begin{lstlisting}
const response = await fetch("https://your-app-domain//items", {
  method: "POST",
  // ...
});
\end{lstlisting}

\item Update README.md file with well-structured information to document your API. It should include:
    
    
\begin{itemize}
\item Brief description of your application.

\item Technologies used (Node.js, Django, Asp.NET).

\item Step-by-step setup instructions, prerequisites for local setup and how to run project locally.
\end{itemize}


\end{enumerate}